\section{Discussion}
\subsection{Answer to the research question}
On this 4-core / 8-thread laptop, the OpenMP version of the blocked kernel scales almost linearly up to the number of \emph{physical} cores
and then stops: $S(4)=3.40$--$3.95$ ($E(4)=0.85$--$0.99$) for every tested size, while 6, 8 or 16 threads change the speedup by at most a few
percent for $n\ge500$. In the main campaign parallel execution paid off at every size, down to $n=100$ (0.31~ms per sequential product, $S(4)=3.54$); the
supplementary experiment located the threshold between $n=16$ and $n=32$: at $n=16$ (1.5~$\mu$s of work) every parallel configuration is slower
than the baseline, at $n=32$ (10~$\mu$s) four threads are 2.6$\times$ faster. The practical answer is therefore: parallelise once a product
costs more than roughly ten microseconds of sequential work, and use as many threads as physical cores; more threads cost efficiency without a
reliable gain.

\subsection{Hypotheses against the evidence}
\begin{itemize}
\item \textbf{H1 supported.} $E(4)\ge0.85$ for all $n$, above the predicted 0.7, and $S(4)\le3.95$ for every size. Two observations
support the frequency explanation without proving it (no frequency was measured): runs with $p\ge3$ slow down by 4--8\% during a process at
$n=3000$--$4000$ while single-thread runs do not (Table~\ref{tab:drift}), as expected if the clock falls under sustained all-core load (the chip's
maximum turbo clock, 3.6~GHz, is well above its 2.7~GHz base \cite{nbc}); and the slowdown appears although \texttt{CPU\_Speed\_Limit} stays at
100, i.e.\ without OS-level throttling. Load imbalance cannot explain
the gap: rows are split to within one row and every thread receives rows.
\item \textbf{H2 supported.} $S(8)/S(4)=0.85$--$1.06$ and $E(8)\le0.52<E(4)$. The vectorisation campaign adds evidence for the mechanism
\cite{marr}: going from 4 to 8 threads (Hyper-Threading) raised the throughput of the scalar build at $n=1000$ by 23\% (median time
from 235 to 192~ms), whose single thread leaves execution units idle, but not that of the SSE2 build (+2\%) and slightly lowered that of the AVX2 build ($-5\%$), whose single thread already keeps the
floating-point units of a core busy.
\item \textbf{H3 partly supported.} Oversubscription never gave a consistent gain over $p=8$ ($S(16)/S(8)=0.74$--$1.31$, with the high
value inside a 24--28\% IQR) and clearly hurt only at $n=100$ ($-26\%$), where each thread's share (6--7 rows, about 20~$\mu$s of work) is
comparable to the cost of time-slicing 16 threads on 8 contexts and waiting at the barrier for the last one. For larger $n$ the cost was
not distinguishable from the noise.
\item \textbf{H4 partly rejected.} The existence of a threshold was confirmed, but its position was wrong: the prediction of $S(4)<2$ for $n\le250$
failed ($S(4)=3.54$ and $3.56$ at $n=100$ and $250$), and the threshold lies between $n=16$ and $n=32$ (supplementary experiment, added after
the main results). The OpenMP runtime keeps its thread pool between parallel regions, so after the warm-up each region only has to restart the
existing threads and synchronise at the final barrier. The estimated cost of a region, 1.4--3.5~$\mu$s for 2--8 threads, grows with $p$ and is of
the same order as the whole sequential product at $n=16$ (1.47~$\mu$s), which explains why more threads are \emph{worse} there; at $n=32$ the
10~$\mu$s of work already amortise it. This conclusion is bound to the timing boundary: products were timed in batches after a warm-up; a
single cold call, which includes the creation of the threads, was not measured separately.
\item \textbf{H5 supported within the noise.} The one-thread parallel run differs from the baseline by $-3\%$ to $+2\%$ for six sizes; the two
exceptions ($-6$ to $-7\%$) go in the \emph{opposite} direction of an overhead and coincide with the noisiest sequential runs. Any OpenMP overhead
with one thread is smaller than the between-process variability ($\approx\pm5$--$10\%$, Section~\ref{sec:var}).
\item \textbf{H6 supported.} Peak RSS = $3\cdot8n^2$ + 0.7--0.8~MiB (sequential) + 0.4--0.6~MiB with OpenMP, independent of $p$ within 0.34~MiB.
Each thread only allocates a stack and runtime bookkeeping; the matrices are shared, not copied.
\item \textbf{H7 partly supported.} AVX2 $>$ SSE2 $>$ scalar at every $p$ and $n$, as predicted. The SIMD gain shrank with threads only for AVX2 at
$n=2000$ (4.40 to 4.01); at $n=1000$ it was not monotonic. The gains exceed the vector widths (2.2--2.9 for width 2, 3.9--5.1 for width 4). A
plausible explanation, consistent with the compiler remarks, is that the vectoriser also interleaves the loop (2 or 4 independent vectors per
iteration), so loop control, address arithmetic and load instructions are spread over 4--16 multiply--adds instead of one; without hardware
counters this remains an explanation, not a measurement.
\end{itemize}

\subsection{Why the speedup saturates at the physical cores}
With FMA disabled, each update $c_{ij}\mathrel{+}=a_{ik}b_{kj}$ needs one multiply and one add. The blocked kernel keeps its tiles in the
private L1/L2 caches \cite{lam,goto}, so it is limited by those floating-point units and by load bandwidth, not by main memory: the
sequential rate is almost constant (5.7--6.6~GFLOP/s) from $n=100$ (all three matrices fit in cache) to $n=4000$ (367~MiB), and the parallel rate is also flat
(21--24~GFLOP/s with $p=4$). In roofline terms \cite{roofline} the kernel stays under the compute roof for every size, which is why the speedup does
not depend on $n$ and why adding hardware threads that share those units (H2) or software threads that share the contexts (H3) cannot help.
The tile-size check points the same way: $S(4)$ was 3.54--3.86 for every tile from 32 to 192.

\subsection{Limitations}
\begin{itemize}
\item \textbf{One machine and one OS.} The results describe one Intel Skylake laptop under macOS; another CPU (more cores, Apple Silicon with
performance and efficiency cores, different turbo policy) can change the shape of the curves.
\item \textbf{No affinity, no counters, no frequency.} Threads could not be pinned, so their placement on cores and hardware threads was chosen by the
scheduler and may differ between runs. Frequency, cache misses and port usage were not measured (\texttt{powermetrics} needs administrator rights);
statements about turbo, caches and execution units are supported indirectly.
\item \textbf{Between-process and between-period variability.} The same configuration differed by up to 17\% between the two passes, much more than the
within-process IQRs. Absolute times should be read with about $\pm10\%$ uncertainty; speedups, whose two terms were measured minutes apart at the
same size, were more stable (largest change 0.36 for $p=16$). The cause of the slower first part of the campaign is unknown. With 5 samples per
configuration (3 at $n=4000$) quartiles are coarse.
\item \textbf{Timing boundary.} Batches after a warm-up measure steady-state cost per product; the cost of the first parallel region (thread creation)
is only present in the warm-up samples and is not analysed.
\item \textbf{Tile size.} $b=96$ was kept for comparability with Assignment~2, although $b=64$ (sequential) and $b=32$ (4 threads) were 5--11\% faster here.
\end{itemize}

\section{Conclusions}
Within the tested scope (blocked i-k-j kernel, $b=96$, \texttt{float64}, $n=100$--$4000$, one i7-6820HQ laptop under macOS, steady-state timing):
\begin{itemize}
\item Row-partitioned OpenMP parallelisation of the Assignment~2 kernel is correct and race-free: all 125 measured products (110 main + 15 supplementary) and 1404 test products are
bit-identical to the sequential baseline, and the reported thread counts match the real OS threads.
\item The speedup grows almost linearly up to the 4 physical cores ($S(4)=3.4$--$4.0$, $E(4)=0.85$--$0.99$) and saturates there; Hyper-Threading
($p=8$) and oversubscription ($p=16$) add no reliable gain for this compute-bound kernel and halve or quarter the efficiency.
\item Parallel execution stops paying off between $n=16$ and $n=32$ (about 1.5 vs.\ 10~$\mu$s of sequential work per product, in steady state):
at $n=16$ every thread count is slower than the baseline and more threads are worse; from $n=32$ upwards 2--4 threads give $S=1.7$--$4.0$.
\item The OpenMP version with one thread has no overhead measurable above the $\pm5$--$10\%$ variability between processes, and memory does not grow with the
number of threads.
\item (Extension) SIMD code generation and threads multiply: AVX2 is 3.9--5.1$\times$ faster than scalar code at equal thread count, and AVX2 with
4 threads is $14.9\times$ faster than the scalar sequential version at $n=2000$ (43.9~GFLOP/s).
\end{itemize}
\textbf{Unresolved:} the cost of a single cold parallel call (thread creation), the exact break-even between $n=16$ and $32$, the cause of the
period-to-period variability, and a direct measurement of clock frequency and cache behaviour, which would need hardware counters and administrator access.

\section{Use of generative AI}
In line with the course policy, the use of AI is disclosed: the tool was \textbf{Claude Code} (Anthropic). It was used for planning the experimental
design and hypotheses, writing the C implementation (adapting the Assignment~2 kernel), the benchmark driver, the test, analysis, plotting and
vectorisation-evidence scripts, the README, a first draft of this report and the oral-defence notes. It also searched for and checked the bibliographic
data of the new references. The author approved the plan and the design decisions, supervised the execution of the measurement campaign on the author's own laptop, and
reviewed the report and the oral-defence guide. The author remains responsible for understanding the code, validating the measurements and defending
the conclusions.

\section*{Reproducibility materials}
Repository: \repourl{} (code, Makefile, README with exact commands, raw CSV files, environment record, analysis scripts and figures). Reused material:
the blocked kernel, input generator, test and driver structure of \cite{a2} (commit \texttt{1aab86d}) and the i-j-k kernel and inputs of \cite{a1}
(commit \texttt{101af1c}); the inputs are regenerated and verified byte for byte with the SHA-256 digests in \texttt{data/manifest.json}.

\begin{thebibliography}{99}\small
\bibitem{a1} H.~Gonz\'alez Fern\'andez, ``Assignment 1: How runtime and memory scale in the basic triple-loop matrix multiplication in C, Java and Python,'' Big Data, Univ. Las Palmas de Gran Canaria, Oct. 2026. Code: \url{https://github.com/hugogzfz/bd-assignment1-matmul}, commit 101af1c.
\bibitem{a2} H.~Gonz\'alez Fern\'andez, ``Assignment 2: Optimised dense matrix multiplication and sparse matrices in C,'' Big Data, Univ. Las Palmas de Gran Canaria, Oct. 2026. Code: \url{https://github.com/hugogzfz/bd-assignment2-matmul}, commit 1aab86d.
\bibitem{lam} M.~D.~Lam, E.~E.~Rothberg and M.~E.~Wolf, ``The cache performance and optimizations of blocked algorithms,'' in \emph{Proc. 4th Int. Conf. Architectural Support for Programming Languages and Operating Systems (ASPLOS IV)}, Santa Clara, CA, USA, 1991, pp.~63--74.
\bibitem{goto} K.~Goto and R.~A.~van de Geijn, ``Anatomy of high-performance matrix multiplication,'' \emph{ACM Trans. Math. Softw.}, vol.~34, no.~3, pp.~1--25, 2008.
\bibitem{openmp} OpenMP Architecture Review Board, \emph{OpenMP Application Programming Interface, Version 5.2}, Nov. 2021. [Online]. Available: \url{https://www.openmp.org/specifications/}
\bibitem{amdahl} G.~M.~Amdahl, ``Validity of the single processor approach to achieving large scale computing capabilities,'' in \emph{Proc. AFIPS Spring Joint Computer Conf.}, Atlantic City, NJ, USA, Apr. 1967, pp.~483--485, doi: 10.1145/1465482.1465560.
\bibitem{marr} D.~T.~Marr, F.~Binns, D.~L.~Hill, G.~Hinton, D.~A.~Koufaty, J.~A.~Miller and M.~Upton, ``Hyper-Threading technology architecture and microarchitecture,'' \emph{Intel Technology Journal}, vol.~6, no.~1, pp.~4--15, Feb. 2002.
\bibitem{roofline} S.~Williams, A.~Waterman and D.~Patterson, ``Roofline: An insightful visual performance model for multicore architectures,'' \emph{Commun. ACM}, vol.~52, no.~4, pp.~65--76, Apr. 2009, doi: 10.1145/1498765.1498785.
\bibitem{nbc} ``Intel Core i7-6820HQ notebook processor,'' Notebookcheck.net. [Online]. Available: \url{https://www.notebookcheck.net/Intel-Core-i7-6820HQ-Notebook-Processor.149416.0.html}. Accessed: Nov.~12, 2026.
\bibitem{llvmvec} LLVM Project, ``Auto-vectorization in LLVM,'' LLVM documentation. [Online]. Available: \url{https://llvm.org/docs/Vectorizers.html}. Accessed: Nov.~12, 2026.
\bibitem{higham} N.~J.~Higham, \emph{Accuracy and Stability of Numerical Algorithms}, 2nd~ed. Philadelphia, PA, USA: SIAM, 2002.
\bibitem{hoefler} T.~Hoefler and R.~Belli, ``Scientific benchmarking of parallel computing systems: Twelve ways to tell the masses when reporting performance results,'' in \emph{Proc. Int. Conf. High Performance Computing, Networking, Storage and Analysis (SC '15)}, Austin, TX, USA, 2015, Art.~no.~73.
\bibitem{numpy} C.~R.~Harris \emph{et al.}, ``Array programming with NumPy,'' \emph{Nature}, vol.~585, pp.~357--362, 2020, doi: 10.1038/s41586-020-2649-2.
\end{thebibliography}
